% Generated draft: requires mathematical and source review.
\documentclass[11pt]{article}
\usepackage[margin=25mm]{geometry}
\usepackage{fontspec}
\setmainfont{Latin Modern Roman}
\usepackage{amsmath,amssymb,mathtools,unicode-math}
\setmathfont{Latin Modern Math}
\usepackage{enumitem,needspace,xurl,hyperref}
\setlength{\parindent}{0pt}
\setlength{\parskip}{5pt}
\setlength{\emergencystretch}{3em}
\newcommand{\N}{\mathbb{N}}
\newcommand{\Z}{\mathbb{Z}}
\newcommand{\Q}{\mathbb{Q}}
\newcommand{\R}{\mathbb{R}}
\newcommand{\C}{\mathbb{C}}
\newcommand{\F}{\mathbb{F}}
\newcommand{\A}{\mathbb{A}}
\newcommand{\PP}{\mathbb{P}}
\newcommand{\E}{\mathbb{E}}
\newcommand{\eps}{\varepsilon}
\newcommand{\dd}{\,\mathrm{d}}
\begin{document}
\section*{UnsolvedMath ID: 700619}
\textbf{Source problem code:} catalogue entry
\subsection{Definitions and mathematical statement}
Let \(G\) be a group. Expand the usual group signature by a set of nullary operation symbols \((c_a)_{a\in A}\), interpreted as specified elements of \(G\). A \emph{mixed identity} is an equation \(t(x_1,\ldots,x_n)=s(x_1,\ldots,x_n)\) between terms in this expanded signature that holds for every assignment \((x_1,\ldots,x_n)\in G^n\). Thus the added symbols act as constants, and the variables are universally quantified.

The requested construction is a class \(\mathcal K\) of groups which is not a mixed variety in the sense intended by the cited work of V. S. Anashin, and which satisfies the following closure requirements:
\begin{enumerate}
\item For every \(G\in\mathcal K\) and every normal subgroup \(N\trianglelefteq G\), the quotient \(G/N\) belongs to \(\mathcal K\).
\item For every nonempty set \(I\) and every family \((G_i)_{i\in I}\) of groups in \(\mathcal K\), the unrestricted Cartesian product \(\prod_{i\in I}G_i\) belongs to \(\mathcal K\).
\item For every \(G\in\mathcal K\), every nonempty set \(I\), and every subgroup \(H\leq G^I\),
\[
\Delta_I(G)\leq H\quad\Longrightarrow\quad H\in\mathcal K,
\qquad
\Delta_I(G)=\{(g)_{i\in I}:g\in G\}.
\]
Here \(G^I=\prod_{i\in I}G\), with coordinatewise group operations.
\end{enumerate}
The source does not specify conventions for empty products or empty powers. If \(\mathcal K\) is nonempty, the quotient condition already implies that the trivial group lies in \(\mathcal K\), so these empty cases add no further closure requirement.

The source defines mixed identities but does not give a definition of a \emph{mixed variety}. In particular, it does not specify how constant symbols and their interpretations are organized across different groups, or how identities determine the class of underlying groups. Consequently, the requirement that \(\mathcal K\) not be a mixed variety remains dependent on the cited theory; a fully self-contained formal statement of that requirement cannot be recovered from the supplied record alone.
\subsection{Short English statement}
Construct a class of groups closed under quotients, Cartesian products, and subgroups of Cartesian powers containing the diagonal subgroup, but which is not a mixed variety in Anashin's intended sense. The supplied definition of mixed identity uses named group elements as constants; the definition of mixed variety still requires clarification from the cited theory.
\subsection{Sources}
\begin{itemize}
\item Supplied local UnsolvedMath registry record, permanent ID \(700619\), exact source problem code \texttt{catalogue entry}, titled \emph{Kourovka Notebook Problem 10.2}. The registry labels the problem \emph{open}; this is supplied source metadata, not an independently verified assessment.
\item \emph{The Kourovka Notebook (Unsolved Problems in Group Theory)}, original problem 10.2, proposed by V. S. Anashin in the 10th issue (1986). The supplied record identifies the 21st issue edition (2026), \texttt{arXiv:1401.0300v48}. The collection description names E. I. Khukhro and V. D. Mazurov as editors.
\item V. S. Anashin, \emph{Math. USSR Sbornik}, volume 57 (1987), pp.\ 171--182, cited in the supplied statement as an example of work on mixed varieties of groups. No article title or further bibliographic details are supplied. The cited text has not been independently consulted for this draft.
\end{itemize}
\end{document}
